
\subsection*{Remarque 3.8.}

Dans le cas @@M0@@ noethérien et @@M1@@ régulier, 3.7 détermine entièrement le groupe abélien @@M2@@ à partir des @@M3@@-anneaux @@M4@@, @@M5@@, des applications @@M6@@ entre ces anneaux, et de l'élément @@M7@@ de @@M8@@ : en effet, @@M9@@ peut se calculer en fonction de @@M10@@ et de @@M11@@ d'après (VI 1.1). Nous verrons au paragraphe suivant que la structure multiplicative de @@M12@@ est déterminée par les mêmes données, de même que la @@M13@@-structure si l'on tensorise par @@M14@@ (cette restriction étant sans doute inutile : voir 4.5).

\section*{4. Immersions régulières et filtration du @@M15@@}

On garde dans ce paragraphe les notations et les hypothèses de 3.1.

\textbf{Lemme 4.1.} \uline{Supposons qu'il existe sur} @@M16@@ \uline{un faisceau localement libre de type fini} @@M17@@, \uline{tel que}, \uline{si} @@M18@@ \uline{sont les classes dans} @@M19@@ \uline{de} @@M20@@, \uline{on ait} @@M21@@. \uline{Alors la congruence suivante est vérifiée} :
@@M22@@

Soit @@M23@@ le rang de @@M24@@ (qu'on peut supposer constant) ; celui de @@M25@@ est alors @@M26@@. Par définition :
@@M27@@
Or @@M28@@. Donc pour tout @@M29@@, @@M30@@ est divisible par @@M31@@. Par suite :
@@M32@@
